The 3D comparisons use the Richardson extrapolation of the 0D block at ne and 2ne radial
elements (ne = 32 for the passive cases, 16 for the active case). To check this reference
without relying on extrapolation, the 0D block was run with up to ne = 256 (passive) and
128 (active) elements (\texttt{zero\_d\_convergence.py}). Table~\ref{tab:zero-d} compares every
direct solution and every extrapolation with the best estimate of the converged 0D
solution, the extrapolation of the two finest levels. The direct solutions converge at
second order and the extrapolations at fourth. The passive 3D reference (32/64) is within
$8\cdot10^{-9}$ of the best estimate in the volume, $4\cdot10^{-8}$ in $\rho$ (full
Holzapfel-Ogden wall) and $3\cdot10^{-7}$ in the twist rate (helical wall; the 0D state does
not depend on the tube length). The direct
solution with ne = 256 is much less accurate (volume up to $2.7\cdot10^{-7}$, $\rho$
$1.3\cdot10^{-6}$, twist rate $3.7\cdot10^{-5}$): it would be good enough for nr $\le 8$, but
not for the full wall at nr = 16, whose volume difference is $7.6\cdot10^{-8}$. The
active reference (16/32) is within $2\cdot10^{-8}$ in the volume and $7\cdot10^{-8}$ in the
pressure. In the very thick wall ($R_e/R_i = 2.5$) the direct 0D solutions have larger
errors (volume $4\cdot10^{-5}$ at ne = 32 instead of $1.7\cdot10^{-5}$), yet the 32/64 reference is
within $8\cdot10^{-8}$ in the volume and $4\cdot10^{-7}$ in $\rho$; for the thin wall and the
large deformation (best estimate 64/128) it is within $6\cdot10^{-10}$. The extrapolated
references are therefore far below every 3D difference reported here.
